\documentclass[11pt,a4paper]{article}
\usepackage{ftm}

\title{Yang–Mills e Gap de Massa\\[4pt]
       \large Caderno de Lacunas Formais — Programa FTM}
\author{Federated Trivalent Manifold Research Programme}
\date{2026}

\begin{document}
\maketitle
\tableofcontents
\bigskip

%% ---------------------------------------------------------------
\section{Enunciado oficial}
%% ---------------------------------------------------------------

\begin{definition}[Yang–Mills e gap de massa]
Seja $G$ um grupo de gauge compacto simples (e.g.\ $SU(N)$).
O problema pede:
\begin{enumerate}[leftmargin=2em]
  \item Construir rigorosamente uma teoria quântica de Yang–Mills
    não trivial em $\mathbb R^4$ satisfazendo axiomas convenientes
    (Osterwalder–Schrader ou Wightman).
  \item Provar que a teoria tem gap de massa $\Delta>0$:
    o espectro do Hamiltoniano é contido em $\{0\}\cup[\Delta,\infty)$,
    com $\Delta$ independente do volume e da regularização.
\end{enumerate}
\end{definition}

\bigskip

%% ---------------------------------------------------------------
\section{Programas de pesquisa e encaixe FTM}
%% ---------------------------------------------------------------

\begin{programs}

\item \textbf{QFT construtiva — programa de Balaban.}
Balaban (1982–1988) desenvolveu um programa de renormalização
em malha para Yang–Mills em 3D, controlando o fluxo de
renormalização via estimativas uniformes sobre blocos de escala.
A extensão para 4D continua incompleta: o controle das
contribuições de alta frequência em 4D exige estimativas
mais refinadas que em 3D.

\begin{gap}
A extensão do programa de Balaban para $\mathbb R^4$
(em vez de $\mathbb R^3$) não está completa; em particular,
a renormalização a todas as ordens com controle uniforme
do limite contínuo não foi provada. \OPEN
\end{gap}

\item \textbf{Confinamento e Wilson loop.}
A lei de área para o Wilson loop,
\[
  \langle W_C\rangle \sim \exp(-\sigma\,\mathrm{Area}(C)),
  \quad \sigma>0,
\]
é esperada em Yang–Mills conformal e é equivalente ao confinamento
de quarks. A teoria do dual-supercondutor (Mandelstam–'t Hooft)
propõe condensação de monopolos como mecanismo.

\begin{gap}
Nenhuma prova rigorosa da lei de área para $SU(N)$ em 4D contínuo. \OPEN
\end{gap}

\item \textbf{Gap espectral e instantons.}
A análise semiclássica de instantons sugere gap de massa
via efeitos não perturbativos de ordem $e^{-8\pi^2/g^2}$,
mas o controle perturbativo diverge e o controle não perturbativo
exige construção rigorosa da integral de caminho.

\item \textbf{Axiomas de Jaffe–Witten e existência da medida.}
O problema do Milênio pede, formalmente, uma medida em espaço
de conexões gauge-equivalentes satisfazendo os axiomas
de Osterwalder–Schrader, com gap positivo e não trivialidade.

\end{programs}

%% ---------------------------------------------------------------
\section{Objeto FTM de gauge}
%% ---------------------------------------------------------------

\begin{definition}[FTM de gauge em escala $a$]
Para grupo compacto simples $G$ e escala de malha $a>0$,
o objeto FTM é
\[
  F_a = (K_a,\Gamma_a,P_a,A_a,S_a),
\]
onde:
\begin{itemize}[leftmargin=2em]
  \item $K_a$ é a malha cúbica ou trivalente de passo $a$;
  \item $P_a$ é um $G$-fibrado sobre $K_a$;
  \item $A_a$ é uma conexão (holonomias por aresta);
  \item $S_a=\sum_{\text{plaquetas}} \mathrm{Re}\operatorname{Tr}(1-U_\square)$
    é a ação de Wilson.
\end{itemize}
As transições entre cartas satisfazem descida de fibrados e
covariância de gauge (cociclo $G$-valorado).
\end{definition}

%% ---------------------------------------------------------------
\section{Obrigações formais FTM para Yang–Mills}
%% ---------------------------------------------------------------

\begin{obligation}[YM1 — Regularização]\label{ob:YM1}
Para cada $a>0$ e volume finito $\Lambda$, definir medida $\mu_{a,\Lambda}$
no espaço de conexões discretas com positividade de reflexão
(ou condição de positividade equivalente).
\end{obligation}

\begin{obligation}[YM2 — Limites uniformes]\label{ob:YM2}
Provar estimativas para funções de correlação
$\langle\mathcal O_1\cdots\mathcal O_k\rangle_{a,\Lambda}$
independentes de $a$, do volume $\Lambda$ e da escolha de atlas FTM.
\end{obligation}

\begin{obligation}[YM3 — Limite contínuo]\label{ob:YM3}
Provar existência de limite $a\to0$, $\Lambda\to\mathbb R^4$,
resultando em distribuições temperadas (funções de Schwinger)
satisfazendo os axiomas de Osterwalder–Schrader, independentes
da regularização escolhida.
\end{obligation}

\begin{obligation}[YM4 — Não trivialidade]\label{ob:YM4}
Provar que o limite contínuo não é a teoria livre ($g=0$)
nem uma teoria degenerada.
\end{obligation}

\begin{obligation}[YM5 — Gap de massa]\label{ob:YM5}
Provar que para observáveis gauge-invariantes,
\[
  \langle\mathcal O(x)\,\mathcal O(0)\rangle_c
  \;\le\; C\,e^{-\Delta|x|},
  \qquad \Delta>0,
\]
com $\Delta$ sobrevivendo aos limites $a\to0$ e $\Lambda\to\infty$.
\end{obligation}

\begin{remark}
Um espectro discreto num grafo finito não prova gap de massa:
é necessário que $\Delta$ seja uniforme e sobreviva a ambos os limites.
A descida FTM é naturalmente adaptada à colagem de fibrados e conexões;
o núcleo construtivo quântico (YM1–YM5) permanece \OPEN.
\end{remark}

\MATURITY{M1 conceitual. YM1–YM5 todas ABERTAS.}

\end{document}
